\section{Methods}
\label{sec:methods}

\subsection{CT decomposition and projection}
\paragraph{Data} We used the LIDC-IDRI collection (1\,018 chest CT series from 1\,010 patients) with the radiologists' nodule contours.
Series were read from DICOM or from the LUNA16 MetaImage copy (voxel-identical on all 12 series we compared, including nodule slice locations); 1\,000 series
could be rendered (18 were skipped because of non-axial orientation or irregular slice spacing). A render was \emph{usable} if the lungs were not cut by the
scan, at most 50 of the 512 frame rows lay beyond the scan, the lung volume was between 1.5 and 9~L and the bone volume exceeded 300~cm$^3$
(844 usable CTs). Patients were assigned to train/validation/test (70/10/20\%) by a hash of the patient identifier (588/98/158 CTs; 1\,057/172/298
consensus nodules) before any model was trained; the split file hash is recorded in \path{freeze.json}.

\paragraph{Components} Each CT was split into bone, non-lung soft tissue and lung. The body is the largest component of $\mathrm{HU}>-300$
(slice-wise hole filling, then a $1{\times}9{\times}9$ opening to detach the couch). Bone follows the iterative-relaxation recipe of the reference method \citep[Algorithm~1]{angaitkar2026anatomy}, with a 10~mm support cube,
the first pass starting from the reference mask and a single closing after the ladder:
a reference mask at 300~HU, thresholds $\tau\in\{200,150,100\}$~HU admitted only where at least 5\% of a 10~mm neighbourhood was already bone, union with the
reference mask, a radius-2 closing to fill marrow, and removal of components smaller than 1\% of the largest. Lung is air ($<-320$~HU) inside the body in
components at least 20\% the size of the largest, closed and hole-filled slice-wise (so vessels and nodules belong to the lung component) and
without bone. Bone voxels were replaced in the soft-tissue component by a normalised convolution of the surrounding soft tissue (Gaussian, 4~mm, with
12 and 36~mm fall-backs). Density is water-equivalent, $\rho=\mathrm{clip}(1+\mathrm{HU}/1000,0,3)$~g\,cm$^{-3}$.

\paragraph{Projection} Each component was integrated along the anteroposterior axis (parallel beam), giving an areal density $a_c$ (g\,cm$^{-2}$).
The frame is a body crop: width 1.08 times the body width, height the lung extent plus 60~mm clipped to 0.75--1.2 of the width (short scans cannot
reach 1.2), area-resized to $512\times512$ and later squashed to a square exactly as real radiographs are. Padded rows are set to air.
The AP view is the horizontal mirror of the PA view.

\paragraph{Detector physics} For a component $c$ with areal density $a_c$, a tube spectrum $\phi(E)$ and an energy-integrating detector,
a pixel value is
\begin{equation}
P_c=-\ln\frac{\sum_E \phi(E)\,E\,\exp\!\big(-(\mu/\rho)_c(E)\,a_c\big)}{\sum_E\phi(E)\,E},
\label{eq:physics}
\end{equation}
with ICRU-44 \citep{icru1989report44} cortical bone, soft tissue and lung attenuation (table~4 of NIST SRD 126; \citealp{hubbell2004xraymass}) and W-anode spectra (SpekPy, 12$^\circ$, 2.5~mm Al) at
70--140~kVp (training draws 80--130~kVp, the synthetic evaluation uses 100~kVp). Components are therefore stored as areal densities and the
spectrum is applied at training time, so one render serves every tube voltage.
The sum of separately projected components differs from the projection of the whole body because of beam hardening; we report this
additivity gap (6.1\% at 100~kVp for 20~g\,cm$^{-2}$ of soft tissue and 3~g\,cm$^{-2}$ of bone). Fig.~\ref{fig:pipeline} shows the
decomposition and projection of one CT.

\begin{figure}[htbp]
\centering
\safefig[0.86\linewidth]{figures/fig_preprocessing_pipeline.pdf}
\caption{CT decomposition and polychromatic projection for one LIDC-IDRI test CT. (a) CT slice; (b) bone, lung and non-lung soft-tissue
masks; (c)--(e) areal densities of bone, non-lung soft tissue and lung; (f) composite radiograph $X$ at 100~kVp; (g) exact bone image $B$;
(h) exact soft-tissue image $S=X-B$; (i) a real radiograph (SZCH-X-Rays) for comparison.}
\label{fig:pipeline}
\end{figure}

\subsection{Training images with exact components}
A training sample is composed on the GPU from the stored areal densities in this order: lesion insertion (below), one random affine transform shared by all maps
(rotation $\pm15^\circ$, scale 0.9--1.15, shift $\pm6\%$, mirror), the physics above, component weighting and a monotone tone curve
$\gamma$:
\begin{equation}
\begin{gathered}
S=\mathrm{norm}\big(P_{\mathrm{nl}}+g\,P_{\mathrm{lung}}\big),\quad B=\frac{w_b\,P_{\mathrm{bone}}}{p_{99}(P_{\mathrm{bone}})},\quad
X=\frac{S+B}{1.5},\\
X'=\gamma(X),\quad S'=\gamma(S),\quad B'=X'-S',
\end{gathered}
\label{eq:compose}
\end{equation}
with $g\sim U(1,3)$ and $w_b\sim U(0.2,0.5)$. The tone curve is either a random gamma or histogram matching to the intensity quantiles of a
randomly drawn real radiograph (4\,000 NIH ChestX-ray14 training-list images). Defining $B'=X'-S'$ keeps the decomposition exact under any
monotone mapping. Gaussian blur is linear and applied identically to $X'$ and $B'$; quantum noise and scatter
modify only $X'$.

\subsection{Suppressor}
\begin{figure}[htbp]
\centering
\safefig{figures/fig_architecture.pdf}
\caption{Model architectures. (a) Study I: bone suppressor trained with the twin lesion-invariance regulariser (Section~\ref{sec:method_v2});
the twin $x^{+}=\min(x+n,1)$ shares the weights and the bone target of $x$. (b) Study II: multi-stream classifier with a shared encoder and a
class-query decoder (Section~\ref{sec:classifier}). (c) Study III: 2.5D Swin CT teacher distilled into a Swin X-ray student on paired
projections (Section~\ref{sec:xm}, Eq.~\eqref{eq:kd}).}
\label{fig:architecture}
\end{figure}
The suppressor (Fig.~\ref{fig:architecture}a) is a SwinV2-Tiny encoder (ImageNet-1k, window 8, $512^2$ input) with a nested (UNet++) decoder over the four Swin stages, CBAM attention, a plain
U-Net tail with two shallow convolutional skips to full resolution and a $1\times1$ head that outputs the bone image directly; the suppressed image is
$x-\hat b$. Relative to the reference this is a smaller backbone at half the resolution, forced by compute. Training used AdamW (lr $2\times10^{-4}$,
weight decay 0.01, 500-step warm-up, cosine decay), batch 8, mixed precision, gradient clipping at 1, EMA weights and 200\,000 samples per arm.
For a predicted bone image $\hat b=f(x)$ and the target $b$, the suppression loss is
\begin{equation}
\begin{split}
\mathcal{L}_{\mathrm{sup}}(\hat b,b)={}&0.5\,\lVert\hat b-b\rVert_1+0.5\,\lVert\hat b-b\rVert_2^2+0.5\,\big(1-\mathrm{SSIM}(\hat b,b)\big)\\
&+2.5\,\big(\lVert\hat b-b\rVert_1+\lVert\hat b-b\rVert_2^2\big)_{\mathrm{ROI}},
\end{split}
\label{eq:loss}
\end{equation}
with norms taken as means over pixels and the ROI taken from the lung component (top structures: $P_{\mathrm{lung}}/p_{99.9}>0.3$). Every arm changes exactly one factor of the recipe arm (Table~\ref{tab:arms}); all other settings,
including the seed set, are identical.

\subsection{Lesion-faithful supervision}
Each lesion (LIDC consensus of at least two readers, 50\% voxel agreement) was projected as its excess density over the surrounding lung.
\emph{Lesion-weighted loss} adds $(\mathrm{L1}+\mathrm{MSE})$ on a dilated lesion mask with weight 1. \emph{Counterfactual lesion insertion (CLIA)}
adds, before the physics, 0--3 lesion patches drawn from \emph{training} CTs at random lung positions (including over bone) with random scale
(0.7--1.5) and an amplitude of 0.6--1.6 times that scale; the bone target is untouched, so the network is explicitly taught that a lesion on a rib is not bone.

\subsection{Evaluation on real images}
\paragraph{Paired dual-energy data} The SZCH-X-Rays test set has 75 radiographs with dual-energy-subtraction soft-tissue references.
Tone differs between a radiograph and its DES image, so metrics are reported as is (``raw'') and after a per-image least-squares gain and
offset of the estimate $\hat s$ onto the reference $s$ (``calibrated'', primary):
\begin{equation}
\begin{gathered}
(\hat a,\hat c)=\arg\min_{a,c}\sum_p\big(a\,\hat s_p+c-s_p\big)^2,\\
\mathrm{PSNR}_{\mathrm{cal}}=10\log_{10}\frac{1}{\frac{1}{|\Omega|}\sum_{p\in\Omega}\big(\mathrm{clip}_{[0,1]}(\hat a\,\hat s_p+\hat c)-s_p\big)^2},
\end{gathered}
\label{eq:psnrcal}
\end{equation}
with intensities in $[0,1]$ and $\Omega$ the evaluated region; SSIM$_{\mathrm{cal}}$ uses the same calibrated estimate. Calibrated PSNR and
SSIM are also reported inside the region where DES removed the most signal (top 10\% of calibrated radiograph minus DES, the bone region)
and outside it.
\paragraph{Rib benchmark} On VinDr-RibCXR (245 images with per-rib masks; evaluated locally, never used for training) we compare the
gradient magnitude $G(I)=\lvert\nabla(g_{1}*I)\rvert$ (Sobel operator after a Gaussian blur of 1~px) of the suppressed image $\hat s$ with that
of the radiograph $x$:
\begin{equation}
\begin{gathered}
R_{\mathrm{rib}}=1-\frac{\sum_{p\in\mathcal{R}}G(\hat s)_p}{\sum_{p\in\mathcal{R}}G(x)_p},\\
R_{\mathrm{non}}=1-\frac{\sum_{p\in\mathcal{N}}G(\hat s)_p}{\sum_{p\in\mathcal{N}}G(x)_p},
\end{gathered}
\label{eq:rib}
\end{equation}
where $\mathcal{R}$ is the rib-boundary band inside the lungs and $\mathcal{N}$ the lung pixels away from ribs. We report both, their
difference $R_{\mathrm{rib}}-R_{\mathrm{non}}$ (selectivity) and the ROC-AUC of the gradient-reduction map for separating the two sets of
pixels.
\paragraph{Controls and baselines} Identity, CLAHE, a homomorphic filter and the publicly released CT-supervised weights of the reference method
(research licence) are scored with the same code.

\subsection{Statistics}
All comparisons are paired. Confidence intervals are percentile bootstrap intervals (10\,000 resamples) over patients when several units come from one patient;
tests are Wilcoxon signed-rank tests with Holm correction within a hypothesis family; non-inferiority uses the lower bound of a one-sided 95\% interval against a
margin fixed in the analysis plan; AUCs are compared with the paired DeLong test. Seeds are averaged within unit before testing and their spread is reported.
The analysis plan, the split and the code commit were frozen before any locked evaluation set was scored (\path{freeze.json}).
